\begin{solapartat}
\[ {}^{18}_{9}F \rightarrow {}^{A}_{B}Y + {}^{C}_{D}\textit{positró} + {}^{0}_{0}\nu \]

\punts{0,2} El nucli de fluor té 9 protons segons veiem en l'equació. El nombre de neutrons seran: $A - Z = 18 - 9 = 9$ neutrons

\punts{0,2} El positró és l'antipartícula de l'electró per tant $C = 0$ i $D = 1$\\
Com que en la desintegració s'ha de conservar el nombre atòmic i el màssic,
\begin{align*}
18 &= A + 0 + 0 \Rightarrow A = 18\\
9 &= B + 1 + 0 \Rightarrow B = 8
\end{align*}

El positró i l'electró s'anihilen donant lloc a 2 fotons idèntics que viatjaran en la mateixa direcció i sentit contrari. L'energia dels 2 fotons serà la que emmagatzemava la massa en repòs de les dues partícules que s'anihilen.

\punts{0,2} $E = 2h\nu = 2mc^2 \Rightarrow h\nu = mc^2 \Rightarrow \nu = \dfrac{mc^2}{h}$

\punts{0,4} $\nu = \dfrac{9,11\times 10^{-31}\times\left(3,00\times 10^{8}\right)^2}{6,63\times 10^{-34}} = 1,24\times 10^{20}\ \mathrm{Hz}$
\end{solapartat}

\begin{solapartat}
\punts{0,2} $N_\mathit{final} = 0,01N_0$

\punts{0,2} $0,01N_0 = N_0 e^{-\lambda t}$

\punts{0,4}
\[ \left.\begin{aligned}
&\ln 0,01 = -\lambda t \Rightarrow t = -\frac{\ln 0,01}{\lambda}\\
&0,5N_0 = N_0 e^{-\lambda\cdot 109,77'} \Rightarrow \lambda = -\frac{\ln 0,5}{109,77'}
\end{aligned}\right\}\ t = 729,30\ \mathrm{min} \approx 12\ \mathrm{h} \qquad \text{\punts{0,2}} \]
\end{solapartat}
